\section{Dalla lettura della issue alla creazione del branch}

\subsection{Step 1: Identificazione della issue}

Lo sviluppatore identifica sulla repository GitHub (\url{https://github.com/Idrorepellente/quantsys_institutional}) una issue etichettata con la label \texttt{strategy}, che richiede l'implementazione di una nuova strategia.

Una issue di richiesta strategia ben formata contiene:

\begin{itemize}
  \item \textbf{Titolo}: nome della strategia (es. \emph{Implementare Statistical Arbitrage Cointegration su S\&P 100});
  \item \textbf{Razionale}: motivazione economica e/o letteratura di riferimento;
  \item \textbf{Edge atteso}: Sharpe target unleveraged, decorrelazione attesa rispetto al resto del sistema;
  \item \textbf{Criteri di accettazione}: soglia di Sharpe OOS, MaxDD massimo, eventuale finestra temporale di validazione richiesta;
  \item \textbf{Riferimenti}: paper accademici, articoli, discussioni precedenti.
\end{itemize}

\subsection{Step 2: Lettura della documentazione tecnica}

Prima di iniziare lo sviluppo, è obbligatorio leggere o rileggere il documento \texttt{GUIDE\_STRATEGIE\_CUSTOM.md} nella radice della repository. Tale documento descrive:

\begin{itemize}
  \item Struttura attesa della cartella della strategia;
  \item Campi obbligatori e opzionali del \texttt{StrategyDescriptor};
  \item Contract minimo che l'implementation deve rispettare;
  \item Sistema di indicator pool condiviso;
  \item Sistema di exit policy (TP/SL);
  \item Procedura di test e walk-forward;
  \item Esempi reali delle strategie già implementate.
\end{itemize}

\subsection{Step 3: Aggiornamento del branch \texttt{main} locale}

Posizionarsi sul branch \texttt{main} locale e aggiornarlo con le ultime modifiche dal remoto:

\begin{lstlisting}
git checkout main
git pull origin main
\end{lstlisting}

\subsection{Step 4: Creazione del branch dedicato}

La convenzione per i branch di strategia è:

\begin{lstlisting}
git checkout -b feature/<nome-strategia>
\end{lstlisting}

dove \texttt{<nome-strategia>} è uguale al nome della cartella che ospiterà la strategia all'interno di \texttt{strategies/}. Esempio: per implementare la strategia \texttt{stat\_arb\_coint}:

\begin{lstlisting}
git checkout -b feature/stat-arb-coint
\end{lstlisting}

Si noti che, per convenzione, il nome del branch usa il trattino come separatore mentre il nome della cartella usa l'underscore.
